\section{Canonical source and parameters}\label{sec:notation}

For an integer \(Y\ge2\), let
\[
 \Omega_Y=\{(m,n)\in\mathbb Z_{>0}^2:
       m<n\le Y,\ Y/2<n,\
       \operatorname{rad}(m(m+1))=\operatorname{rad}(n(n+1))\}.
\]
An element of \(\Omega_Y\) is called a source pair.
For each source pair define its four corner gcds by
\[
A=\gcd(m,n),\quad B=\gcd(m,n+1),\quad
C=\gcd(m+1,n),\quad D=\gcd(m+1,n+1).
\]
These four integers are pairwise coprime. For example, a prime dividing
\(A\) and \(B\) would divide both \(n\) and \(n+1\). A prime
dividing \(A\) and \(C\) would divide both \(m\) and \(m+1\). The
same two consecutive pairs settle the other four pairs. Each corner
divides the two entries in its definition. Hence, by pairwise
coprimality, \(AB\mid m\), \(CD\mid m+1\), \(AC\mid n\), and \(BD\mid n+1\).
Define the positive integer residuals
\[
a=\frac m{AB},\qquad b=\frac{m+1}{CD},\qquad
c=\frac n{AC},\qquad d=\frac{n+1}{BD}.
\]
\begin{samepage}
The residuals are pairwise coprime. First, \(\gcd(a,b)=\gcd(c,d)=1\),
since each pair divides consecutive integers. For the four cross pairs,
divide the defining corner gcds. This gives
\[
\gcd(Ba,Cc)=\gcd(Aa,Dd)=\gcd(Db,Ac)=\gcd(Cb,Bd)=1,
\]
so \(\gcd(a,c)=\gcd(a,d)=\gcd(b,c)=\gcd(b,d)=1\) as well.
\end{samepage}
The coordinates and mask are
\[
\begin{gathered}
m=ABa,\quad m+1=CDb,\quad n=ACc,\quad n+1=BDd,\\
Z=ABCD,\quad R=\operatorname{rad}(Z),\quad M=abcd,\\
K_{\rm res}=M/\operatorname{rad}(M),\quad L=cd-ab,\\
\alpha=ac,\quad\beta=bd,\quad\gamma=ad,\quad\delta=bc,\\
u=(m+n+1)/(BC),\qquad t=(n-m)/(AD),\\
CDb-ABa=1,\qquad BDd-ACc=1.
\end{gathered}\AutoTag{eq:old:2_1}
\]
The four cross-gcd identities above will be used below. The radical
condition is exactly \(\operatorname{rad}(M)\mid Z\). If \(p\mid a\), then
\(p\mid m\). Equality of the original radicals forces \(p\mid n\) or
\(p\mid n+1\). Hence \(p\mid A\) or \(p\mid B\). The other residuals give
the corresponding corner. Conversely, every prime of \(Z\) divides both
consecutive products. The mask ensures that no residual introduces a
prime outside \(Z\).
Thus the common radical is
\[
\operatorname{rad}(m(m+1))=\operatorname{rad}(n(n+1))=R.
\]
The identities
\[
Z^2M=m(m+1)n(n+1),\quad ZL=(n-m)(m+n+1)
\]
give \(ut=L>0\) once the shifts are integral. Their integrality follows
from the corner identities
\[
m+n+1=C(Ac+Db)=B(Aa+Dd),\qquad
n-m=A(Cc-Ba)=D(Bd-Cb).
\]
Since \(\gcd(B,C)=\gcd(A,D)=1\), the first numerator is divisible by
\(BC\) and the second by \(AD\). Thus \(u,t\) are positive integers
with product \(L\).
The shifts are units at the opposite corner pairs:
\[
\gcd(u,AD)=\gcd(t,BC)=1.
\]
Modulo a prime dividing \(A\), the integer \(m+n+1\) is
congruent to \(1\). Modulo a prime dividing \(D\), it is congruent to
\(-1\). Modulo a prime dividing \(B\), \(n-m\) is congruent to
\(-1\). Modulo a prime dividing \(C\), it is congruent to \(1\).
The denominator factors \(BC\) and \(AD\) are
coprime to these opposite corner pairs.

The two conics, scale identities, and inverse maps from the same
parametrization are
\[
\begin{aligned}
\beta D^2-\alpha A^2&=u,&
\delta C^2-\gamma B^2&=t,\\
\alpha A^2&=u\frac{mn}{m+n+1},&
\beta D^2&=u\frac{(m+1)(n+1)}{m+n+1},\\
\gamma B^2&=t\frac{m(n+1)}{n-m},&
\delta C^2&=t\frac{n(m+1)}{n-m},\\
B&=(cA+bD)/u,& C&=(aA+dD)/u,\\
A&=(dB-bC)/t,& D&=(cC-aB)/t.
\end{aligned}\AutoTag{eq:old:2_2}
\]
We use only source points. They satisfy positivity, integrality, the
four corner gcds and the cross gcds. They also satisfy the shell
condition and the radical mask. We assert no converse for arbitrary
positive solutions of either conic.

The selected chart is deterministic:
\[
(J,K;s)=
\begin{cases}(A,D;u),&u\le t,\\(B,C;t),&t<u.\end{cases}\AutoTag{eq:old:2_3}
\]
Thus \(u=t\) belongs to the first chart. Write \((f,g)=(\alpha,\beta)\)
in the first chart and \((f,g)=(\gamma,\delta)\) in the second. For a
positive integer \(x\), let \(v_p(x)\) be the exponent of \(p\) in its
prime factorization. For a finite set of primes \(\mathcal P\), put
\[
x_{\mathcal P}=\prod_{p\in\mathcal P}p^{v_p(x)},\qquad
x_{\varnothing}=1.
\]
Then
\[
f_J=f_{\{p:p\mid J\}},\quad g_K=g_{\{p:p\mid K\}},\quad
v=f/f_J,\quad w=g/g_K,\quad
r_J=\operatorname{rad}(f_J),\quad r_K=\operatorname{rad}(g_K),
\]
\[
\Pi=\min(r_J,v)\min(r_K,w).\AutoTag{eq:old:2_4}
\]
These supported parts are functions of the canonical source. They are
not extra summation variables.

Set
\[
\begin{gathered}
S_0=\frac{30869}{5000000},\quad
\epsilon_0=k_0=S_0,\quad l_0=2S_0,\quad j_0=p_0=3S_0,\\
a_0=\frac6{125},\quad b_0=\frac{4166537}{30000000},\quad
\Delta_{\rm exp}=3S_0,\\
q_{\min}=\frac{9907393}{15000000}=\frac23-S_0,\quad
q_{\max}=\frac{5092607}{7500000}=\frac23+2S_0,\quad
\tau_0=1-\frac32S_0 .
\end{gathered}\AutoTag{eq:old:2_5}
\]
\begin{samepage}
We use the following notation.
\begin{center}
\small
\begin{tabular}{@{}ll@{}}
Symbol & Meaning\\
\hline
\(Y\) & dyadic shell endpoint; \(Y/2<n\le Y\)\\
\(A,B,C,D\) & four corner gcds\\
\(a,b,c,d\) & positive residual factors\\
\(Z,R,M,L\) & \(ABCD,\ \operatorname{rad}(Z),\ abcd,\ cd-ab\)\\
\(\alpha,\beta,\gamma,\delta\) & \(ac,\ bd,\ ad,\ bc\)\\
\(u,t\) & positive shifts, with \(ut=L\)\\
\(J,K;s\) & selected corner pair and its smaller shift\\
\(f_J,g_K,v,w,\Pi\) & supported coefficients, complements, and their mask\\
\(S_0,g_0,\beta_0\) & minimum saving, reserve, and allocated loss (values in the Introduction)\\
\end{tabular}
\end{center}
\end{samepage}
Define the high sector
\[
\mathcal C_{\rm hi}=\{(m,n)\in\Omega_Y:
R>Y^{5/3-\epsilon_0},\
M>Y^{q_{\min}},\
\operatorname{rad}(M)>Y^{2/3-k_0}\}.
\AutoTag{eq:old:2_6}
\]
The high sector is the complement of the first three predicates in
\eqref{eq:old:3_1}. Those rows test the complementary inequalities.
For sources with \(R>Y^{5/3-\epsilon_0}\), the source identities give
\[
M<\kappa_Y^2Y^{q_{\max}},\qquad
L<\kappa_YY^{1/3+\epsilon_0},\qquad \kappa_Y=1+Y^{-1}.
\AutoTag{eq:old:2_7}
\]
We have \(m(m+1)n(n+1)\le\kappa_Y^2Y^4\) and
\(ZL=(n-m)(m+n+1)<\kappa_YY^2\). Also
\(Z>Y^{5/3-\epsilon_0}\). On \(\mathcal C_{\rm hi}\), the strict lower
cuts on \(R\) and \(M\) also give
\[
Y^{4-3\epsilon_0}<Z^2M=m(m+1)n(n+1)
                       \le2\kappa_Ym^2Y^2,
\]
so \(m^2>Y^{2-3\epsilon_0}/(2\kappa_Y)
\ge Y^{2-3\epsilon_0}/3\). Therefore
\[
m>Y^{\tau_0}/\sqrt3.\AutoTag{eq:old:2_8}
\]
Here \(2\kappa_Y\le3\) for \(Y\ge2\). This last inequality does not
use the radical cut in \(\mathcal C_{\rm hi}\).

After the \(Q0\) row of \eqref{eq:old:3_1} has failed,
\(K_{\rm res}=M/\operatorname{rad}(M)
\le Y^{\Delta_{\rm exp}}\). After \(Q1\) has failed,
\(L\ge Y^{1/3-l_0}\). After both \(Q2\) chart tests have failed, the
selected corners satisfy \(J,K\ge Y^{1/3-j_0}\). These are the
non-strict complements used in the following partition.
Display~\eqref{eq:old:3_1} states the predicates and their order.
